\section{Methods}
\subsection{Dataset and evaluation protocol}
The benchmark iterates over all images stored in \path{datasets/images} and matches each image with a binary ground-truth mask in \path{datasets/groundthruth} using the naming rule \path{<image_stem>_mask.png}. In the present repository state, the dataset contains four histological TIFF images: \path{ROI_1.tiff}, \path{ROI_2.tiff}, \path{ROI_3.tiff}, and \path{ROI_4.tiff}. Three images have spatial resolution \(4001 \times 4001\) pixels and one image has resolution \(4000 \times 4000\) pixels. Ground-truth masks are binary PNG files with values \(\{0,255\}\).

All methods expose a unified interface,
\begin{verbatim}
run_model(image_path) -> np.ndarray
\end{verbatim}
which returns a binary mask of type \texttt{uint8}. The benchmark resizes each predicted mask to the corresponding ground-truth resolution using nearest-neighbor interpolation to ensure pixelwise comparability. Reproducibility is promoted by fixed random seeds (\texttt{seed=42}) for Python, NumPy, and PyTorch where applicable.

\subsection{VeSpA pipeline}
The VeSpA implementation is imported from \path{pipelines/VeSpA/Segmentation + Measurements.py}. The pipeline reads each color image in BGR format, converts it to a CMYK representation, and uses the yellow channel as the main contrast source for vessel extraction. The yellow channel is smoothed by Gaussian filtering and binarized with Otsu thresholding. The resulting binary image is refined through morphological closing, dilation, and erosion, followed by hole filling and cleanup operations using closing and opening. The final binary vessel mask is returned as the segmentation output. This pipeline is deterministic and does not rely on learned weights.

\subsection{SAM pipeline}
The SAM implementation is provided in \path{pipelines/sam/run_sam.py}. The current benchmark uses the ViT-B SAM checkpoint stored in \path{models/sam_vit_b_01ec64.pth}. Input images are converted from BGR to RGB and downscaled to a maximum side length of 1024 pixels before inference in order to control memory use and execution time. Automatic mask generation is then applied using the repository configuration \path{points_per_side=16}, \path{pred_iou_thresh=0.8}, \path{stability_score_thresh=0.9}, \path{crop_n_layers=0}, and \path{min_mask_region_area=64}. Candidate masks covering more than 25\% of the resized image area are filtered out to suppress large background regions. The remaining masks are merged by union and resized back to the original image resolution.

\subsection{YOLOv8 segmentation pipeline}
The YOLO implementation is provided in \path{pipelines/yolo/run_yolo.py}. The benchmark loads the pretrained Ultralytics segmentation model from \path{models/yolov8n-seg.pt}. Each image is processed with inference size \path{imgsz=1024} and \path{retina_masks=True}. When the model returns multiple instance masks, each mask is resized to the original image resolution with nearest-neighbor interpolation and the final binary prediction is obtained as the pixelwise union of all instances.

\subsection{Evaluation metrics}
Four binary segmentation metrics are implemented in \path{evaluation/metrics.py}: Dice coefficient, intersection over union (IoU), precision, and recall. Let \(P\) denote the predicted foreground set and \(G\) the ground-truth foreground set. The metrics are computed as
\[
\mathrm{Dice} = \frac{2|P \cap G|}{2|P \cap G| + |P \setminus G| + |G \setminus P|},
\]
\[
\mathrm{IoU} = \frac{|P \cap G|}{|P \cap G| + |P \setminus G| + |G \setminus P|},
\]
\[
\mathrm{Precision} = \frac{|P \cap G|}{|P \cap G| + |P \setminus G|},
\qquad
\mathrm{Recall} = \frac{|P \cap G|}{|P \cap G| + |G \setminus P|}.
\]
When the denominator of a metric is zero, the implementation returns 1.0.

\subsection{Outputs generated by the benchmark}
The benchmark writes per-image and per-model metrics to \path{results/metrics/results.csv}, model-level means to \path{results/metrics/summary.csv}, and mean runtime per image to \path{results/metrics/runtime_summary.csv}. For qualitative inspection, the pipeline also saves binary predictions and color overlays in \path{results/prediction/}. Summary visualizations consist of a Dice boxplot and a mean-IoU bar chart saved under \path{results/metrics/}.
